\section*{Capítulo \ref{ch:g12:destruction-of-mountains} --- A destruição das montanhas}

\begin{solution}{exo:g12:destruction-of-mountains:1}
(a) a \omterm{def:g12:destruction-of-mountains:denudation}{taxa de denudação}; (b) a \omterm{def:g12:destruction-of-mountains:load}{carga de fundo}; (c) a \omterm{def:g12:destruction-of-mountains:load}{carga dissolvida};
(d) a \omterm{def:g12:destruction-of-mountains:diagenesis}{diagênese}; (e) o \omterm{def:g12:destruction-of-mountains:rebound}{rebote isostático}.
\end{solution}

\begin{solution}{exo:g12:destruction-of-mountains:2}
(1) Um conglomerado. (2) Um arenito. (3) Um folhelho. (4) Uma brecha.
\end{solution}

\begin{solution}{exo:g12:destruction-of-mountains:3}
\qty{1}{m/Myr} são \qty{1000}{mm} em \num{1000000} anos, \qty{0.001}{mm}
por ano. (1) \qty{0.054}{mm} por ano. (2) \qty{0.218}{mm} por ano.
(3) $\qty{10}{mm}$ por ano são \qty{10}{m} em 1000 anos, \qty{10000}{m/Myr}.
\end{solution}

\begin{solution}{exo:g12:destruction-of-mountains:4}
A carga sólida é \qty{85}{\percent} do total, logo o total é $60/0.85
\approx 70.6$ milhões de toneladas, e a \omterm{def:g12:destruction-of-mountains:load}{carga dissolvida} cerca de $70.6 - 60 \approx
10.6$ milhões de toneladas por ano.
\end{solution}

\begin{solution}{exo:g12:destruction-of-mountains:5}
Intemperismo do granito, transporte por um rio, deposição numa bacia,
soterramento sob camadas mais jovens, cimentação.
\end{solution}

\begin{solution}{exo:g12:destruction-of-mountains:6}
A \qty{30}{cm/s}, o grão de areia é levantado (acima de \qty{20}{cm/s}); o
grão de silte continua em suspensão. A \qty{5}{cm/s}, o grão de areia se
deposita (abaixo de \qty{10}{cm/s}), mas o grão de silte continua em
suspensão, já que \qty{0.1}{cm/s} basta para mantê-lo ali.
\end{solution}

\begin{solution}{exo:g12:destruction-of-mountains:7}
Carga total $4/0.85 \approx 4.71$ milhões de toneladas, $\num{4.71e9}/2830 \approx
\qty{1.66e6}{m^3}$ de rocha por ano. Sobre $\qty{5e10}{m^2}$: $\num{3.3e-5}$ m
por ano, cerca de \qty{33}{m/Myr}, abaixo da mediana de \qty{54}{m/Myr}.
\end{solution}

\begin{solution}{exo:g12:destruction-of-mountains:8}
(1) $1 - 2700/3300 \approx 0.182$. (2) Cerca de \qty{182}{m}. (3) Um pouco
mais que os \qty{163}{m} do capítulo: quanto mais leve a crosta em
comparação com o manto, menos a coluna sobe.
\end{solution}

\begin{solution}{exo:g12:destruction-of-mountains:9}
(1) \qty{8}{km} de rocha. (2) $8000 \times 0.163 \approx \qty{1300}{m}$.
(3) $8000/200 = 40$ milhões de anos.
\end{solution}

\begin{solution}{exo:g12:destruction-of-mountains:10}
(1) $794/43 \approx 18$ milhões de anos. (2) A superfície desce apenas
$0.163$ vez a espessura erodida, logo a rocha a remover é $794/0.163
\approx \qty{4900}{m}$, o que leva cerca de $4900/43 \approx 110$ milhões de
anos. (3) Os continentes não são só desgastados: a formação das montanhas e
o magmatismo dos arcos continuam a erguê-los e a acrescentar crosta nova.
\end{solution}

\begin{solution}{exo:g12:destruction-of-mountains:11}
(1) $\num{1e14}/2830 \approx \qty{3.5e10}{m^3}$, isto é, cerca de
\qty{35}{km^3}. (2) $100/15.5 \approx 6.5$ anos.
\end{solution}

\begin{solution}{exo:g12:destruction-of-mountains:12}
$E = kh$ e $h = h_0 e^{-fkt}$ dão $E(t) = kh_0 e^{-fkt}$. A espessura
removida de $t = 0$ a $t \to \infty$ é $\int_0^{\infty} kh_0
e^{-fkt}\,\mathrm{d}t = kh_0/(fk) = h_0/f$. Para abaixar a superfície de
$h_0$, a erosão precisa remover $h_0/f \approx 6.1\,h_0$, como diz a
proposição do rebote: cada metro erodido abaixa a superfície só $f$ metro.
\end{solution}

\begin{solution}{exo:g12:destruction-of-mountains:13}
(1) $10 \times 1000 = \qty{10000}{mm}$, \qty{10}{m}. (2) Remover uma carga
de gelo de espessura $t$ equivale a remover uma massa $\rho_{\text{gelo}}\,t$
por unidade de área; a coluna sobe até deslocar essa mesma quantidade a
menos de manto:
$2000 \times 900/3380 \approx \qty{530}{m}$. (3) A terra sobe cada vez menos
à medida que se aproxima do seu novo equilíbrio: o que impulsiona a subida é
a parte ainda não recuperada.
\end{solution}

\begin{solution}{exo:g12:destruction-of-mountains:14}
(1) Antes de 1900: cerca de \qty{56}{m/Myr} (como no exemplo). Em
1987--2006: $145/0.85 \approx 171$ milhões de toneladas, $\num{1.71e11}/2830
\approx \qty{6.0e7}{m^3}$, sobre $\qty{2.98e12}{m^2}$: cerca de
\qty{20}{m/Myr}. (2) A mais antiga: a carga mais recente é retida por
barragens e outras obras. (3) Não: a diminuição é de $400 - 145 = 255$
milhões de toneladas por ano, das quais as barragens explicam de 100 a 150,
cerca da metade.
\end{solution}

\begin{solution}{exo:g12:destruction-of-mountains:15}
(1) Pelo menos \qty{6}{km}. (2) $310 - 252 = 58$ milhões de anos no máximo,
logo pelo menos $6000/58 \approx \qty{100}{m/Myr}$. (3) Fica entre a mediana
(\qty{54}{m/Myr}) e a média (\qty{218}{m/Myr}) das bacias: uma taxa comum
para um cinturão de montanhas.
\end{solution}

\begin{solution}{pb:g12:destruction-of-mountains:1}
\textbf{1.} $E = 0.209 \times 2000 \approx \qty{418}{m/Myr}$.

\textbf{2.} \qty{0.418}{mm} por ano.

\textbf{3.} Mais rápida que ambas: é uma cadeia íngreme e jovem, tão rápida
quanto as bacias sismicamente ativas, e mais.

\textbf{4.} A sua declividade média.

\textbf{5.} Uma camada de $\qty{0.418e-3}{m}$ sobre $\qty{1e6}{m^2}$ são
\qty{418}{m^3}, $418 \times 2830 \approx \qty{1.18e6}{kg}$: cerca de 1200
toneladas.

\textbf{6.} $2000/418 \approx 4.8$ milhões de anos.

\textbf{7.} $h(t) = 2000\,e^{-0.209\,t}$ ($t$ em milhões de anos).

\textbf{8.} $1/0.209 \approx 4.8$ milhões de anos.

\textbf{9.} $\ln 2/0.209 \approx 3.3$ milhões de anos.

\textbf{10.} $e^{-0.209 \times 11} = e^{-2.30} \approx 0.100$.

\textbf{11.} $1 - 2830/3380 \approx 0.163$.

\textbf{12.} A superfície desce $100 \times 0.163 \approx \qty{16}{m}$; a
\omterm{def:g12:continental-crust:root}{raiz} sobe cerca de \qty{84}{m}.

\textbf{13.} $h' = -0.163 \times 0.209\,h \approx -0.0341\,h$, logo $h(t) =
2000\,e^{-0.0341\,t}$.

\textbf{14.} $(2000 - 200)/0.163 \approx \qty{11000}{m}$ de rocha.

\textbf{15.} Rochas que estavam a cerca de \qty{11}{km} de profundidade:
rochas metamórficas como xistos e gnaisses, e granitos.

\textbf{16.} $1/0.0341 \approx 29$ milhões de anos.

\textbf{17.} $29 \times \ln 10 \approx 67$ milhões de anos.

\textbf{18.} $2000\,e^{-0.0341 \times 200} \approx 2000 \times 0.0011 \approx
\qty{2}{m}$.

\textbf{19.} $2000\,e^{-0.0341 \times 250} \approx 2000 \times 0.0002 \approx
\qty{0.4}{m}$: uma planície plana. Um cume de \qty{2037}{m} significa que o
modelo sozinho não basta: a região deve ter sido erguida de novo desde
então, ou erodida mais lentamente do que o modelo supõe.

\textbf{20.} A altitude média se reduz à metade em $29 \times \ln 2 \approx$
\textbf{20 milhões de anos}.
\end{solution}
